% ============================================================
% V. VALIDATION METRICS AND TUNING
% ============================================================
\section{Validation Metrics and Tuning}
\label{sec:metrics}

To allow users to evaluate and compare approximation quality quantitatively, the toolchain defines a set of per-link and aggregate metrics.
These metrics serve both as diagnostic tools during parameter tuning and as pass/fail gates in automated pipelines.

% ---- A. Coverage Metrics ----
\subsection{Coverage}

\paragraph{Coverage boolean.}
The most fundamental requirement is that every mesh vertex lies within at least one capsule.
Coverage is satisfied when $\max_i \mathrm{signed\_distance}_i \leq \epsilon$, where $\mathrm{signed\_distance}_i$ is the signed distance from vertex $\mathbf{v}_i$ to its nearest capsule (negative when inside).

\paragraph{Worst signed distance.}
When coverage is not achieved, the worst (most positive) signed distance identifies the maximum uncovered gap.
This metric is useful for diagnosing which link or region is under-covered.

% ---- B. Volume and Tightness Metrics ----
\subsection{Volume and Tightness}

\paragraph{Capsule union volume (capV).}
The union volume of all capsules in a link is estimated via Monte Carlo sampling.
The bounding box of the capsule set is uniformly sampled on a regular grid with $S$ samples per axis (default $S = 32$, giving approximately $32{,}768$ sample points).
Each sample point is classified as inside or outside the capsule union, and the volume is estimated as:
\begin{equation}
\mathrm{capV} = \mathrm{vol}(\mathrm{AABB}) \cdot \frac{|\{\text{samples inside}\}|}{|\{\text{total samples}\}|}.
\label{eq:capv}
\end{equation}

\paragraph{CapV/AABB ratio.}
The capsule union volume normalized by the link's axis-aligned bounding box volume:
\begin{equation}
\mathrm{capV/aabb} = \frac{\mathrm{capV}}{\mathrm{vol}(\mathrm{AABB})}.
\end{equation}
Values close to 1.0 indicate that the capsule set fills the link's extent efficiently; values significantly above 1.0 suggest excessive over-coverage.

\paragraph{Radius-to-median-bin ratio ($r/\mathrm{binMed}$).}
Along each capsule axis, vertices are binned into ten equal-length segments.
For each bin, the maximum perpendicular distance of assigned vertices to the capsule axis is recorded.
The median of these bin-maximum distances ($\mathrm{binMed}$) characterizes the typical surface offset.
The ratio $r / \mathrm{binMed}$ measures how much the capsule radius is inflated beyond the typical radial extent of its assigned vertices.
High values indicate that a single capsule bridges a tapered or irregular shape that could benefit from splitting.

% ---- C. Primitive Count ----
\subsection{Primitive Count}

The total number of primitives---spheres, capsules, or convex meshes---provides a compactness measure.
For capsule mode, each capsule is counted as one analytic primitive despite being decomposed into three URDF elements (one cylinder and two spheres) during emission.

% ---- D. Configuration Parameters ----
\subsection{Configuration Parameters}

The capsule fitting pipeline exposes a small number of tunable parameters that control the trade-off between primitive count and approximation tightness.
These parameters are grouped by their role in the pipeline.

\paragraph{Sectioning.}
The number of axial slicing planes, denoted $N$ (default $N = 4$), controls the resolution at which cross-sectional geometry is captured.
Higher values produce finer axial sampling at increased computational cost and may yield additional capsules.
Typical values range from 2 (coarse) to 8 (fine).

\paragraph{Circle fitting.}
Per-section circle quality is governed by the circle-outside-area threshold $\tau_{\mathrm{coa}}$ (default $5 \times 10^{-3}$), which bounds the acceptable area of each fitted circle lying outside the contour polygon.
When adaptive circle fitting is enabled, the pipeline incrementally adds circles until the total COA per section falls below $\tau_{\mathrm{coa}}$, up to a per-section maximum $M$ (default $M = 1$).
With adaptive fitting disabled, exactly $M$ circles are fitted per section using deterministic farthest-point initialization and $k$-means clustering.

\paragraph{Capsule budget.}
The per-link capsule budget $K_{\max}$ (default $K_{\max} = 12$) provides a hard upper bound.
When the initial capsule count exceeds $K_{\max}$, the pipeline iteratively removes the capsule whose deletion minimizes the resulting union volume, re-growing remaining capsules to preserve coverage after each removal.

\paragraph{Splitting.}
Capsules that bridge tapered or irregular geometry are detected through two metrics.
The radius-inflation ratio $\rho = r / \mathrm{binMed}$ triggers a split when it exceeds a threshold $\rho_{\max}$ (default 1.45); setting $\rho_{\max} = -1$ disables radius-based splitting.
The volume ratio $\eta = \mathrm{capV} / \mathrm{vol}(\mathrm{AABB})$ triggers a split when it exceeds $\eta_{\max}$; negative values (default $-1$) disable volume-based splitting.
A split is accepted only when the resulting capsule pair reduces the affected union volume by at least $\delta_{\min} = 5 \times 10^{-3}$ in relative terms.

\paragraph{Volume estimation.}
Monte Carlo estimation of capsule union volume uses $S$ samples per axis (default $S = 32$), yielding approximately $S^3$ total sample points.

These parameters define the three named presets---\emph{single} ($N = 2$, $K_{\max} = 1$), \emph{default} ($N = 4$, $K_{\max} = 12$, $\rho_{\max} = 1.45$), and \emph{high\_detail} ($N = 6$, $K_{\max} = 16$), providing sensible starting points for different trade-off regimes.
